% Material migrated from the former appendix into the main body for IPDPS.
% All added prose is review-marked so coauthors can inspect it before finalization.

\subsection{Additional Evaluation Details}
\label{sec:additional-eval-details}

\anurag{Appendix - F Evaluation: Baselines}

\review{\noindent\textbf{Baseline implementations.}
IrEne~\cite{irene} is the closest methodological baseline because it also uses a hierarchical model-tree abstraction. It represents execution at mathematical-operation, machine-learning-operation, and module levels, trains leaf-level energy predictors from runtime resource utilization and model specifications, and recursively combines child predictions to obtain higher-level energy estimates. We use IrEne as designed for single-GPU energy modeling and apply it to the same inference configurations used for \sys{}; its principal limitation in our setting is that communication is not represented as a first-class module.}

\review{CodeCarbon~\cite{codecarbon} is a software-based energy estimator that integrates telemetry-derived power over time, including GPU measurements from NVML and host measurements where available. We include it because it is widely used in practice and represents a strong software-only measurement baseline, but it does not model the structure of parallel communication. For Wilkins et al.~\cite{wilkins2024offlineenergyoptimalllmserving}, we implement the proposed token-in/token-out regression by fitting $e_K(\tau_{\mathrm{in}},\tau_{\mathrm{out}})=\alpha_0\tau_{\mathrm{in}}+\alpha_1\tau_{\mathrm{out}}+\alpha_2\tau_{\mathrm{in}}\tau_{\mathrm{out}}$ on a calibration set. Finally, for NVML we construct a model-level linear regressor whose independent variable is NVML-reported GPU energy and whose target is measured total-system energy. This lets us test whether a readily available GPU-only signal can stand in for a full-system predictor.}

\anurag{Appendix - G Additional results: NVML as a proxy for total energy}

\review{\noindent\textbf{NVML as a total-energy proxy.}
We first ask whether GPU-only energy from NVML can predict full-system energy through a simple regression. On A6000, the direct NVML proxy produces large model-level errors across every family: Vicuna ranges from 29.8--33.4\% MAPE, Mistral from 39.0--44.2\%, Llama from 28.5--31.3\%, and Qwen from 38.0--42.4\%. These errors are substantially larger than \sys{} because NVML omits host-side power as well as communication and synchronization effects that alter the relationship between GPU-board energy and end-to-end system energy.}

\review{We also test whether the NVML regression generalizes to an unseen model variant. In a leave-one-variant-out setting, the NVML-based predictor incurs 47.4--57.3\% MAPE, averaging 51.5\% across the evaluated models, whereas \sys{} is in the 15.8--24.5\% range on the same A6000 generalization task. The gap grows because a single GPU-energy signal does not encode architectural changes that affect communication volume, synchronization, and non-GPU system energy. This experiment supports the need for the architecture-aware and communication-aware features used by \sys{} rather than treating NVML as a sufficient proxy for total energy.}

\anurag{Appendix - H Additional Results: Generalization across models}

\review{\noindent\textbf{Cross-family generalization.}
Generalizing to an entirely unseen model family is more difficult than holding out one variant because the computation graph, attention mechanism, and module dimensions can all change. We therefore perform leave-one-family-out evaluation: all measurements from one family are excluded during training, and the resulting predictor is tested on that held-out family under the same deployment configurations. On A6000, \sys{} obtains 24.1\%, 27.0\%, 26.1\%, and 27.6\% MAPE when Vicuna, Mistral, Llama, and Qwen are respectively held out, compared with 49.3\%, 56.5\%, 55.3\%, and 58.4\% for IrEne. On B40, \sys{} obtains 13.2\%, 15.1\%, 23.1\%, and 12.2\%, while IrEne incurs 63.8\%, 71.2\%, 113.3\%, and 62.9\%. The result indicates that decomposition into reusable module types plus structural features transfers more effectively across architectures than a single-GPU hierarchy that omits communication.}

\anurag{Appendix - I Additional Results: Generalization across variants}

\review{\noindent\textbf{Cross-variant and workload generalization.}
In leave-one-variant-out evaluation on A6000, \sys{} obtains 15.84\%, 17.72\%, and 17.55\% MAPE for Vicuna-7B/13B/33B; 21.17\%, 23.13\%, and 24.52\% for Mistral-8B/24B/48B; 18.15\%, 21.11\%, and 22.41\% for Llama-7B/13B/70B; and 20.15\%, 20.99\%, and 18.98\% for Qwen-8B/14B/32B. On B40, the corresponding values are 11.83\%, 12.22\%, and 12.81\%; 14.41\%, 15.29\%, and 15.81\%; 13.56\%, 14.64\%, and 23.52\%; and 14.25\%, 17.21\%, and 13.21\%. The larger Llama-70B and Mistral variants are harder than the smaller variants, but the predictor remains substantially more accurate than the baselines.}

\review{We separately test workload extrapolation across batch sizes. Across model families, batch-size generalization produces 16.89--21.82\% MAPE with an average of 19.05\%. Generalizing from batch size 16 to 32 gives 19.33\% average MAPE, while the reverse direction (32 to 16) gives 18.77\%. The similar error in both directions suggests that the execution and structural features capture the scaling relationship between workload size and energy rather than merely memorizing one operating point.}

\anurag{Appendix - J Additional Results: Module-level energy predictions}

\review{\noindent\textbf{Fine-grained module prediction.}
Module-level prediction is important because the purpose of the model tree is not only to improve the final model-level estimate but also to expose where energy is spent. On A6000, for Vicuna, the 2-/4-GPU MAPE is 8.8/11.4\% for Self-Attention, 6.6/9.5\% for MLP, 17.3/19.4\% for AllReduce, 6.4/7.3\% for LayerNorm/RMSNorm, and 9.9/9.6\% for the embedding module. On B40, the corresponding values are 5.3/8.2\%, 4.8/8.8\%, 11.2/13.1\%, 3.2/7.8\%, and 7.1/8.2\%. The communication-heavy AllReduce module is consistently the most difficult component because its energy includes transfer and synchronization variability, but its error remains moderate enough for useful composition.}

\review{Across model families, the same trend holds. On A6000 the average module-level MAPE is 6.28\% for Vicuna, 11.01\% for Mistral, 7.93\% for Llama, and 9.03\% for Qwen; on B40 the corresponding values are 4.1\%, 8.8\%, 4.3\%, and 5.1\%. Mistral is the most challenging family in both cases. Its grouped-query attention and SwiGLU blocks alter the mix of compute and communication inside a Transformer block, which increases the difficulty of module-level prediction. These results nevertheless show that the model-tree abstraction retains accuracy across different module compositions rather than only matching end-to-end totals.}

\anurag{Appendix - K Additional Results: Feature correlation analysis}

\review{\noindent\textbf{Feature correlation analysis.}
To validate that the feature set contains information beyond a single GPU-energy signal, we compute Spearman rank correlations between the candidate features and total measured energy across Vicuna configurations. NVML GPU energy is strongly correlated with total energy, with $\rho=0.633$--$0.762$, which is expected because GPU computation is a major contributor. However, several non-NVML features are at least as informative: batch size has $\rho=0.735$--$0.737$, execution time has $\rho=0.658$--$0.664$, and memory usage has $\rho=0.656$--$0.663$. These correlations show that energy is jointly shaped by workload size, runtime duration, memory activity, and GPU energy. In the multi-GPU case, communication volume and synchronization further perturb that relationship, motivating the multi-feature representation used by \sys{} instead of a single utilization or telemetry scalar.}

\anurag{Appendix - L Additional Results: Using Latency Prediction Approaches to Predict Energy}

\review{\noindent\textbf{Can a latency predictor be retargeted to energy?}
We evaluate whether a latency-prediction methodology can be repurposed by changing only its target label. We implement a NeuSight-style baseline~\cite{neusight}: the execution is decomposed into fine-grained matrix tiles, computational features such as FLOPs and memory behavior are retained, and the same MLP-style prediction structure is used, but the training target is changed from latency to measured system energy. This isolates whether the feature representation that is effective for latency is also sufficient for energy.}

\review{The answer is no. The retargeted predictor obtains 79.19\% MAPE on Vicuna-7B, 93.24\% on Vicuna-13B, and 89.38\% on Vicuna-33B, for an average of 87.3\%. This is roughly 3--5$\times$ worse than \sys{} and 1--2$\times$ worse than the least accurate energy-prediction baseline in our evaluation. Latency-oriented features primarily describe how long kernels execute; they do not explicitly represent power-state behavior, system-level energy, inter-GPU communication volume, or synchronization/waiting effects. The experiment therefore supports the design choice that energy prediction is not simply latency prediction with a different target variable.}

\anurag{Appendix - M Additional Results on NVIDIA H200 GPUs}

\review{\noindent\textbf{Hardware extension to H200.}
We additionally evaluate \sys{} on 2$\times$ NVIDIA H200 GPUs to determine whether the same model-tree decomposition, communication model, and feature set remain useful on a substantially different GPU/interconnect platform. Across the evaluated variants, the model-level MAPE averages 16.8\%. The per-variant errors are 13.34\%, 13.78\%, and 14.45\% for Vicuna-7B/13B/33B; 16.25\%, 17.25\%, and 17.83\% for Mistral-8B/24B/48B; 15.30\%, 16.51\%, and 26.53\% for Llama-7B/13B/70B; and 16.07\%, 19.41\%, and 14.90\% for Qwen-8B/14B/32B. The larger Llama-70B is again the most difficult configuration, but the overall error remains comparable to the A6000 and B40 experiments.}

\review{The H200 module-level results are also low: 8.80\% MAPE for Self-Attention, 7.44\% for MLP, 9.25\% for AllReduce, 8.12\% for LayerNorm/RMSNorm, and 7.78\% for the embedding module. In particular, the single-digit AllReduce error provides evidence that the communication-aware component remains effective when moving from PCIe-oriented platforms to the H200's higher-bandwidth NVLink environment. We do not claim a hardware-agnostic predictor---\sys{} is still calibrated per hardware platform---but these results show that the methodology itself transfers without changing the decomposition or feature design.}
